% 02-abbreviations.tex - Section 2: Abbreviations and Definitions

\chapter{Abbreviations and Definitions}
\label{sec:abbreviations}

\section{Abbreviations}

\begin{longtable}{p{3cm}p{11cm}}
    \toprule
    \textbf{Abbreviation} & \textbf{Definition} \\
    \midrule
    \endfirsthead
    \toprule
    \textbf{Abbreviation} & \textbf{Definition} \\
    \midrule
    \endhead
    \bottomrule
    \endfoot

    APC & American Power Conversion (PDU/UPS manufacturer) \\
    ECS & Environmental Control System \\
    EFM & Electric Field Mill (lightning detector) \\
    EnviroMux & Network Technologies environmental monitoring system \\
    HVAC & Heating, Ventilation, and Air Conditioning \\
    ICD & Interface Control Document \\
    IceQube & IceQube V Series air conditioner \\
    LWA & Long Wavelength Array \\
    M\&C & Monitor and Control \\
    MC4002 & Bard MC4002 HVAC controller \\
    MCS & Monitor and Control System \\
    MIB & Management Information Base \\
    MJD & Modified Julian Date \\
    MPM & Milliseconds Past Midnight \\
    PDU & Power Distribution Unit \\
    SHL & Shelter subsystem \\
    SNMP & Simple Network Management Protocol \\
    UDP & User Datagram Protocol \\
    UPS & Uninterruptible Power Supply \\
    VAC & Volts Alternating Current \\

\end{longtable}

\section{Definitions}

\begin{description}
    \item[Rack] A standard 19-inch equipment rack within the shelter. Each rack may have one or more PDUs.

    \item[PDU] A Power Distribution Unit that provides switched power outlets for equipment. Controlled via SNMP.

    \item[UPS] An Uninterruptible Power Supply that provides battery backup and power conditioning.  Controlled via SNMP.

    \item[Set Point] The target temperature for the HVAC system.

    \item[Differential] The temperature offset from the set point at which cooling engages. For example, with a set point of 72\textdegree F and a differential of 2\textdegree F, cooling starts at 74\textdegree F.

    \item[SCRAM] An emergency shutdown mode that immediately stops all monitoring threads but does not power off equipment.

    \item[Monitor Period] The polling interval for a particular monitoring thread (e.g., temperature, PDU status).
\end{description}
